\documentclass[11pt]{article}
\usepackage[margin=0.75in]{geometry}
\usepackage{setspace}
\usepackage{iftex}
\usepackage{amsmath,amssymb}
\usepackage{graphicx}
\usepackage{enumitem}
\usepackage{pdfpages}
\usepackage{booktabs}
\usepackage{cite}
\usepackage{url}
\usepackage{xcolor}
\usepackage{colortbl}
\usepackage{hyperref}

\ifPDFTeX
  \usepackage[T1]{fontenc}
  \usepackage{newtxtext,newtxmath}
\else
  \usepackage{fontspec}
  \setmainfont{Times New Roman}
\fi

\setstretch{1.15}
\setlength{\parskip}{0.5em}
\setlength{\parindent}{0pt}
\pagenumbering{arabic}

\begin{document}

% ----------------------------------------------------
% PAGE 1: COVER PAGE
% ----------------------------------------------------
\includepdf[pages=1]{doe_p3_cover.pdf}

% ----------------------------------------------------
% PAGE 2: MAIN BODY
% ----------------------------------------------------
\section{Focus Area, Rationale, and Stakeholder Relevance}
This project addresses the strategic siting and capacity sizing of clean energy microgrids and storage systems within power distribution grids. Siting decisions represent a high-stakes, combinatorial planning challenge under complex electrical constraints. To guide real planning outcomes, we model a multi-objective optimization problem that co-optimizes grid resilience ($C_{\rm R}$), capital installation cost ($C_{\rm I}$), and operational voltage control ($C_{\rm VC}$). Our pre-rendered Jupyter notebooks on siting optimization, AI load scaling and $N-1$ resilience are available at~\cite{doe_notebook}.

The primary stakeholders and beneficiaries of this methodology include:
\begin{itemize}[leftmargin=*,noitemsep,topsep=0pt]
  \item \textbf{Electric Utilities and ISOs:} Enables optimal capital allocation for grid assets and helps maintain power quality under dynamic load conditions.
  \item \textbf{Department of Energy (DOE):} Accelerates grid decarbonization and modernization mandates through quantum-enhanced optimization frameworks.
  \item \textbf{Local Communities:} Enhances energy security and resilience against catastrophic outages during extreme weather events.
\end{itemize}

\section{Technical Methodology and Hybrid Quantum Integration}
In our \textbf{first} notebook, we run a hybrid Benders-QAMOO workflow on IEEE-33 bus. Our approach combines Quantum Approximate Multi-Objective Optimization (QAMOO) by~\cite{qamoo2025} with classical Benders decomposition for strategic siting of energy storage and microgrids. Our hybrid workflow is inspired by ~\cite{pasqal2024}, but with a different architecture partition: it divides the optimization into a classical Master Problem (MP) and an AC power-flow Subproblem (SP).
At its core, this is formulated as a mixed-integer linear programming (MILP) optimization problem, which is loaded into a Qiskit QuadraticProgram. We jointly optimize the following objectives inspired by~\cite{multiverse2023}: $C_{\rm I}$: capital investment cost of siting the microgrid assets, $C_{\rm R}$: resilience of the grid, and $C_{\rm VC}$: voltage control quality. This takes the following minimization:
$
\min_{\overline{x} \in \{0,1\}^N} \quad C(\overline{x}) = \lambda_{\rm I} C_{\rm I} + \lambda_{\rm R} C_{\rm R} - \lambda_{\rm VC} C_{\rm VC}
$

In our \textbf{second} notebook, we execute a QAOA pipeline against non-linear structural shocks, stress-testing IEEE 39 bus and, for Phase 3, the larger IEEE 118-bus system (per the challenge's own recommendation for a "tractable" regional-transmission-scale test case) across two severe grid-modernization scenarios: an additive load injection model simulating hyperscale AI data centers, and an exhaustive contingency sweep across transmission lines.

\subsection{Classical Benders Preprocessing (CPU)}
The MP handles the discrete layout coordinates, selecting the binary portfolio of bus locations for asset placement ($\hat{x}^{(t)} \in \{0,1\}^N$). It is solved using CPLEX~\cite{cplex2024}. The candidate layout $\hat{x}^{(t)}$ is passed to the SP evaluator, which runs an AC power-flow simulation using pandapower~\cite{pandapower2018}. If physical line limits (thermal branch loading $> 100\%$) or voltage limits ($0.95 \le V_j \le 1.05$ p.u.) are violated, the SP generates linear feasibility cuts:
\begin{itemize}[leftmargin=*,noitemsep,topsep=0pt]
  \item \textbf{Line Overload Cuts:} For overloaded branches between bus $u$ and $v$: $x_u + x_v \le 1$.
  \item \textbf{Voltage Deviation Cuts:} For bus $w$ experiencing operational bounds violation: $x_w \le 0$.
  \item \textbf{Exact Vector Forbid (Fallback):} If specific buses cannot be isolated: $\sum_{i: \hat{x}_i^{(t)}=1} x_i + \sum_{i: \hat{x}_i^{(t)}=0} (1 - x_i) \le N - 1$.
\end{itemize}
These cuts compile iteratively ($T=5$) to form a set of linear constraints on the master layout solver.

\subsection{QUBO Penalty Mapping}
To embed feasibility constraints into the QAMOO quantum layer, we convert the collected cuts into quadratic penalty terms $p_i = \lambda \cdot \hat{x}_i^{(t)}$ that modify the self-loop diagonals of the target Ising objective graphs, where $\lambda = 10.0$ is the penalty weight. This mutates the diagonal terms of the objective graphs to penalize infeasible portfolios, steering the quantum variational sampler away from voltage and thermal limit violations.

\subsection{Quantum Optimization and Error Mitigation}
Once the penalized graphs are compiled, we train a multi-objective QAOA ($p=3$ layers) optimized using COBYLA, giving us a warm-start on a scalarized representation of the three objectives. The Pareto front is evaluated by executing batched sample runs of size $S=10$ using Dirichlet-distributed scalarization weights.

To overcome physical device noise, we have implemented and integrated the client-side Samplomatic error-mitigation pipeline using Qiskit's new Executor primitive. This pipeline combines three error mitigation techniques: Probabilistic Error Cancellation (PEC), Shaded Lightcones (SLC) and Propagated Noise Absorption (PNA).

\subsection{\texorpdfstring{AI Data Center Shock Modeling and Multi-Objective $N-1$ Contingency Resilience}{AI Data Center Shock Modeling and Multi-Objective N-1 Contingency Resilience}}
Phase 2 validated this mechanism using a synthetic $+300$~MW/$90$~MVAr point-source load injected onto two buses and a uniform $N-1$ sweep across every transmission line. For Phase 3, both inputs are replaced with real, citable data --- see Section~\ref{sec:data-sources} for the AI-load injection model and empirical contingency weighting actually used on the IEEE-118 grid. The underlying QUBO structure is unchanged: cost penalties are modeled with an empirical resilience metric (replacing $C_{\rm I}$ with $C_{\text{resilience}}$), where a per-bus resilience sensitivity index ($R_n$) mutates the self-loop diagonals of the optimization QUBO as follows:
$$Q_{i,i} = -\alpha V_n[i] + \mu r_n[i] - \nu R_n[i] + \lambda(1 - 2K)$$

% =====================================================
% NEW SECTION: Data Sources and Modeling Assumptions
% =====================================================
\section{Data Sources and Modeling Assumptions}
\label{sec:data-sources}
This section documents the provenance of all external datasets used in the Phase~3 pipeline, explicitly delineating what is derived from real-world government data versus what constitutes a modeling assumption.

\subsection{Data Source Provenance}
\begin{table}[h!]
\centering
\caption{External Dataset Provenance}
\label{tab:dataset-provenance}
\vspace{0.5em}
\begin{tabular}{p{3.5cm}p{4.0cm}p{2.7cm}p{3.0cm}}
\toprule
\textbf{Pipeline Input} & \textbf{Source Dataset} & \textbf{DOI / ID} & \textbf{Provenance} \\
\midrule
Total AI load addition ($\Delta P_{\rm DC}$) & PNNL/EPRI IM3+EPRI Data Center Load Projections (TELL BA) & \url{10.57931/3007669} & \cellcolor{green!15}Real govt.\ data \\
\addlinespace
Outage event category probabilities $P(\text{cat})$ & PNNL Event-Correlated Outage Dataset (EAGLE-I $\times$ DOE-417, 2014--2023) & OEDI 6458 & \cellcolor{green!15}Real govt.\ data \\
\addlinespace
Bus-level AI load allocation (which buses receive $\Delta P_{\rm DC}$) & Degree-centrality hub selection with 15\% concentration ratio & --- & \cellcolor{yellow!25}Modeling assumption \\
\addlinespace
Line hazard categorization (Weather / Physical / Equipment) & Thermal loading and degree-centrality percentile thresholds & --- & \cellcolor{yellow!25}Modeling assumption \\
\addlinespace
IEEE bus test systems (14/33/39/118) & pandapower standard library~\cite{pandapower2018} & --- & \cellcolor{blue!10}Synthetic benchmark \\
\bottomrule
\end{tabular}
\end{table}

\subsection{Data Preprocessing Pipeline and Dataset Selection Rationale}
Both real-data inputs above are generated by standalone preprocessing scripts (\texttt{prep\_outage\_scenarios.py}, \texttt{prep\_ai\_load\_overlay.py}) that regenerate model-ready CSVs directly from the cited public sources at setup time, rather than vendoring processed data in the repository --- keeping the codebase lean and every input independently reproducible from its citation.

\textbf{Outage/contingency data.} We selected the PNNL Event-Correlated Outage Dataset (EAGLE-I $\times$ DOE-417, OEDI submission 6458) over a single-year DOE-417 extract because it spans 10 years (2014--2023), which is necessary to compute stable per-category empirical probabilities rather than a noisy single-year estimate. Preprocessing consolidates EAGLE-I's 15-minute county-level outage segments into discrete events by taking the peak customers-affected count per \texttt{(event\_id, county)} pair \emph{before} summing across counties --- summing the raw 15-minute segments directly would double-count outages that persist across multiple reporting intervals within the same county. This yields 663 consolidated events, from which the category probabilities in Table~\ref{tab:outage-probs} are computed as each category's share of total events.

\textbf{AI load data.} We selected the PNNL/EPRI IM3+EPRI TELL dataset over a hand-picked load figure because it provides hourly, Balancing-Authority-level projections across 16 scenarios (4 climate pathways $\times$ 4 EPRI growth rates), letting the load-growth magnitude be cited to a specific public source rather than assumed. Preprocessing filters the dataset to the TVA balancing authority and computes the added data-center demand as the difference between the \texttt{with\_DC} and baseline load columns for the moderate- and high-growth scenarios, identifying the real 2025 peak-total-demand hour directly from the data to anchor the injection scenario.

\subsection{AI Data Center Load Injection Model}
\textbf{What is from real data.}
The total megawatt addition for the TVA Balancing Authority under the high-growth scenario ($\Delta P_{\rm DC} = 383.40$~MW) is sourced directly from the PNNL/EPRI TELL dataset~\cite{pnnl_tell_2025}. This dataset provides hourly Balancing Authority--level load projections with and without data center demand under moderate and high growth scenarios, covering all 66 U.S.\ BAs.

\textbf{What is a modeling assumption.}
The TELL dataset operates at the BA level and does \emph{not} specify which substations within a BA absorb data center demand. Since IEEE benchmark grids are synthetic topologies with no geographic mapping, any sub-BA bus-level allocation is necessarily a modeling choice.

We adopt a \textbf{sparse hub-bus injection model} based on the following empirical evidence of geographic clustering:
\begin{itemize}[leftmargin=*,noitemsep,topsep=0pt]
  \item \textbf{PJM (2025):} Data centers account for 94\% of projected load growth, with demand concentrated in the Northern Virginia corridor (``Data Center Alley'')~\cite{pjm_load_2025}.
  \item \textbf{IEA (2025):} Approximately 50\% of U.S.\ data centers under development cluster near existing large facilities at high-voltage interconnection points~\cite{iea_dc_2025}.
  \item \textbf{LBNL Queued Up (2025):} Approximately 80\% of large-load interconnection requests target fewer than 15\% of available substations~\cite{lbnl_queued_2025}.
  \item \textbf{WECC (2024):} Data center interconnection requests are geographically concentrated at a small number of high-capacity substations~\cite{wecc_assessment_2024}.
\end{itemize}
Based on these sources, we set $f_{\rm hub} = 0.15$ (15\% of load buses) as the concentration fraction. This value is consistent with the LBNL finding that $<$15\% of substations absorb the majority of large-load requests. The selected hub buses are the top-$\lceil f_{\rm hub} \cdot N_{\rm load} \rceil$ buses ranked by graph degree centrality (a proxy for high-voltage interconnection capacity). Within the selected hub set, demand is distributed proportionally to existing bus load:
\begin{equation}
P_n^{\rm shocked} = P_n^{\rm base} + \frac{P_n^{\rm base}}{\sum_{k \in \mathcal{H}} P_k^{\rm base}} \cdot \Delta P_{\rm DC}, \quad n \in \mathcal{H}
\end{equation}
where $\mathcal{H}$ is the set of hub buses and $P_n^{\rm base} = 0$ for $n \notin \mathcal{H}$ (no injection at non-hub buses).

To ensure AC power flow convergence on small test grids (IEEE-14, IEEE-33), we additionally cap the total injection at 25\% of baseline system capacity: $\Delta P_{\rm DC}^{\rm eff} = \min(\Delta P_{\rm DC}, \, 0.25 \cdot \sum_k P_k^{\rm base})$.

\subsection{\texorpdfstring{$N-1$ Contingency Resilience Model}{N-1 Contingency Resilience Model}}
\textbf{What is from real data.}
The empirical event-category probabilities used to weight the $N-1$ contingency metric ($dist_n$) are derived from the PNNL Event-Correlated Outage Dataset~\cite{pnnl_oedi_6458}, which synthesizes 10~years (2014--2023) of EAGLE-I county-level outage records with DOE-417 Electric Emergency Incident and Disturbance Reports, covering 663 major grid disturbance events across the continental United States. The resulting empirical probabilities are:

\begin{table}[h!]
\centering
\caption{Empirical Outage Category Probabilities (OEDI 6458, 2014--2023)}
\label{tab:outage-probs}
\vspace{0.5em}
\begin{tabular}{lcc}
\toprule
\textbf{Event Category} & \textbf{$P(\text{cat})$} & \textbf{Source} \\
\midrule
Weather / Natural Disaster & 44.9\% & \cellcolor{green!15}EAGLE-I $\times$ DOE-417 \\
Physical / Security       & 27.3\% & \cellcolor{green!15}EAGLE-I $\times$ DOE-417 \\
Equipment / Operational   & 24.6\% & \cellcolor{green!15}EAGLE-I $\times$ DOE-417 \\
Other / Unknown           &  3.0\% & \cellcolor{green!15}EAGLE-I $\times$ DOE-417 \\
Cyber                     &  0.2\% & \cellcolor{green!15}EAGLE-I $\times$ DOE-417 \\
\bottomrule
\end{tabular}
\end{table}

\textbf{What is a modeling assumption.}
The mapping from hazard categories to specific transmission lines within IEEE benchmark grids is a modeling assumption, since these synthetic topologies have no geographic or physical-asset metadata. We use a \textbf{physics-grounded percentile categorization} scheme that generalizes across any grid size:
\begin{itemize}[leftmargin=*,noitemsep,topsep=0pt]
  \item \textbf{Weather hazards} $\rightarrow$ Lines in the top-20\% thermal loading percentile (proxy: exposed, high-stress corridors).
  \item \textbf{Physical/Security hazards} $\rightarrow$ Lines touching hub buses in the top-15\% degree centrality percentile (proxy: high-value interconnection targets).
  \item \textbf{Equipment/Operational hazards} $\rightarrow$ Remaining lines (proxy: routine failure modes).
\end{itemize}
For each category, a single representative line (the highest thermally loaded within its category) is taken out of service ($N-1$ contingency), and reactive power probes ($\pm 5.0$~MVAR) at each candidate bus measure voltage recovery improvement ($R_n^{\rm cat}$). The final per-bus resilience metric is:
\begin{equation}
dist_n = \sum_{\text{cat}} P(\text{cat}) \cdot \frac{R_n^{\rm cat}}{\max_k R_k^{\rm cat}}
\end{equation}

\subsection{AC Power Flow Solver Robustness}
Under stressed loading conditions (AI load shock + $N-1$ outage), the standard Newton-Raphson (NR) AC power flow solver may fail to converge. We implement a cascading fallback strategy: NR ($\text{max\_iter}=100$) $\rightarrow$ Iwamoto Damped NR $\rightarrow$ Backward/Forward Sweep (BFSW). This ensures convergence across all test cases without masking genuine physical infeasibility.

\subsection{Summary of Modeling Assumptions}
\begin{table}[h!]
\centering
\caption{Summary of Key Modeling Assumptions and Justifications}
\label{tab:assumptions-summary}
\vspace{0.5em}
\begin{tabular}{p{4.0cm}p{3.5cm}p{5.5cm}}
\toprule
\textbf{Assumption} & \textbf{Value / Choice} & \textbf{Justification} \\
\midrule
Hub concentration fraction $f_{\rm hub}$ & 15\% of load buses & LBNL Queued Up 2025: 80\% of requests at $<$15\% of substations \\
\addlinespace
Hub selection criterion & Degree centrality & Proxy for high-voltage interconnection capacity \\
\addlinespace
Max load shock cap & 25\% of baseline $\Sigma P$ & Power flow convergence on small IEEE grids \\
\addlinespace
Weather line threshold & 80th pctl.\ thermal loading & Exposed, high-stress corridors \\
\addlinespace
Physical/Security bus threshold & 85th pctl.\ degree centrality & High-value interconnection targets \\
\addlinespace
$N-1$ representative line & Max loading per category & Tractable probing (1 line/category vs.\ exhaustive) \\
\bottomrule
\end{tabular}
\end{table}

\section{Data Modeling and Benchmarking Strategy}
We primarily focus on IEEE datasets. Benders-QAMOO is benchmarked against two exact classical methods:
\begin{enumerate}[leftmargin=*,noitemsep,topsep=0pt]
  \item \textbf{CPLEX Weighted-Sum:} Classical solver executing scalarized MILP evaluations.
  \item \textbf{Defining Point Algorithm (DPA):} State-of-the-art MILP algorithm with provable convergence.
\end{enumerate}

\section{Experimental Results and Discussion}
\underline{First Notebook (IEEE-33):} Table~\ref{tab:pareto-ieee33} measures key multi-objective optimization metrics: Pareto front size ($|{\rm PF}|$), Hypervolume (${\rm HV}$), Inverted Generational Distance (${\rm IGD}$), spacing, additive epsilon indicator ($\epsilon^+$), and coverage ($C(A,B)$).

\begin{table}[h!]
\centering
\caption{Pareto Quality Metrics on IEEE 33-Bus Grid}
\label{tab:pareto-ieee33}
\vspace{0.5em}
\begin{tabular}{lccccccc}
\toprule
\textbf{Method} & |\textbf{PF}| & \textbf{HV} & \textbf{IGD} & \textbf{Spacing} & $\epsilon^+$ & \multicolumn{2}{c}{\textbf{Coverage}} \\
\cmidrule(lr){7-8}
& & & & & & $C(A,B)$ & $C(B,A)$ \\
\midrule
QAMOO (Quantum)            & 189 & 707.2790 & 0.0961 & 0.3724 & 0.0588 & 0.5455 & 0.0529 \\
Benders-QAMOO (Quantum)    & 270 & 824.3916 & 0.1097 & 0.2613 & 0.0714 & 0.4697 & 0.0333 \\
CPLEX (Exact Classical WS) & 9   & 47.6514  & 0.1466 & 0.1560 & 0.2569 & 0.1818 & 1.0000 \\
DPA (Exact Classical)      & 57  & 49.1532  & 0.0001 & 0.0665 & 0.0004 & 0.8636 & 1.0000 \\
\bottomrule
\end{tabular}
\end{table}

The empirical metrics reveal a stark architectural trade-off between quantum exploration and classical precision. The quantum models demonstrate massive dominance in exploring non-dominated design spaces, with Benders-QAMOO producing a Pareto front of size $|{\rm PF}| = 270$. This dense coverage yields an order-of-magnitude superiority in Hypervolume (HV), where Benders-QAMOO secures the highest score of 824.3916.

Conversely, the exact classical solvers establish the absolute mathematical baseline for convergence at the expense of frontier diversity. The DPA anchors the true reference front, achieving a near-zero IGD of $0.0001$, a minimal Spacing of $0.0665$, and a near-perfect additive epsilon indicator ($\epsilon^+$) of $0.0004$. This mathematical precision is validated by the asymmetric Coverage indicators, where the exact classical sets achieve a absolute coverage ($C(B,A) = 1.0000$).

\underline{Second Notebook (IEEE-118, Phase 3):} A full Benders-QAMOO run was executed on the real DOE-417/EAGLE-I
contingency-weighted, TVA AI-load-injected IEEE-118 grid. The grid loads correctly at 118 buses; a real
$+383.40$~MW AI load shock (high-growth scenario) is injected across 14 hub buses selected by the top-15\%
degree-centrality rule (buses 11, 14, 16, 31, 48, 53, 55, 58, 69, 74, 76, 79, 91, 99); and the empirical resilience
weighting loads correctly across all 5 contingency categories from the 663 real DOE-417/EAGLE-I events (probabilities
summing to 1.0000, as expected). After pruning to the highest-sensitivity candidates, the QUBO encodes 15 decision
variables (qubits), drawn from candidate buses [19, 20, 21, 27, 42, 43, 50, 51, 52, 57, 107, 108, 113, 114, 116].

\begin{table}[h!]
\centering
\caption{Pareto Quality Metrics on IEEE 118-Bus Grid}
\label{tab:pareto-ieee118}
\vspace{0.5em}
\begin{tabular}{lccccccc}
\toprule
\textbf{Method} & |\textbf{PF}| & \textbf{HV} & \textbf{IGD} & \textbf{Spacing} & $\epsilon^+$ & \multicolumn{2}{c}{\textbf{Coverage}} \\
\cmidrule(lr){7-8}
& & & & & & $C(A,B)$ & $C(B,A)$ \\
\midrule
QAMOO (Quantum)            & 57 & 16.3948 & 0.0000 & 0.1360 & 0.0000 & 1.0000 & 0.0175 \\
Benders-QAMOO (Quantum)    & 76 & 17.3980 & 0.0000 & 0.1043 & 0.0000 & 1.0000 & 0.0132 \\
CPLEX (Exact Classical WS) & 1  & 3.6449  & 0.0000 & N/A    & 0.0000 & 1.0000 & 1.0000 \\
DPA (Exact Classical)      & 0  & 0.0000  & N/A    & N/A    & N/A    & 0.0000 & N/A    \\
\bottomrule
\end{tabular}
\end{table}

On the real IEEE-118 grid, Benders-QAMOO again dominates in front size and Hypervolume, producing $|{\rm PF}| = 76$
points at ${\rm HV} = 17.3980$ (max ${\rm HV} = 281.87$ over the raw non-dominated union), versus $|{\rm PF}| = 57$ and
${\rm HV} = 16.3948$ for standalone QAMOO -- confirming that the classical Benders feasibility cuts continue to sharpen
the quantum search even at this larger 118-bus, 15-qubit scale. CPLEX's exact weighted-sum baseline collapses to a
single scalarized point ($|{\rm PF}| = 1$), which is expected behavior for a weighted-sum method on a non-convex
front and is why $C(B,A) = 1.0000$ (the single CPLEX point is trivially covered by the much larger quantum fronts).
The DPA benchmark could not be run in this environment: the \texttt{dpa-main} binary is unavailable for the qBraid
Linux runtime used for this test (see \texttt{doe/phase\_3/doc/COMPILATION\_GUIDE.md} for build instructions), so its
row is reported honestly as degenerate (${\rm HV}=0$, N/A elsewhere) rather than omitted or fabricated.

\section{Quantum Platform Justification and Resource Needs}
\textbf{Requested:} 156-qubit Heron-class QPUs (e.g., Kingston, Fez) (Inferior alternative: 120-qubit Eagle-class).\\
Our rationale for selecting Heron over Eagle:
\begin{itemize}[leftmargin=*,noitemsep,topsep=0pt]
  \item \textbf{Scale:} Heron supports up to 156 qubits, enabling direct variable encoding of grid cases up to 156 bus widths.
  \item \textbf{Fidelity:} Heron's high 2-qubit gate fidelity ($> 99\%$) and low cross-talk are required for Samplomatic to resolve the signal (fails on noisier Eagle systems).
  \item \textbf{Integration:} Native support for Qiskit Executor primitives.
\end{itemize}

\section{Gap Analysis and Implementation Roadmap}
The following tasks are part of our Phase 3 roadmap:
\begin{itemize}[leftmargin=*,noitemsep,topsep=0pt]
  \item \textbf{True Quantum Hybrid Benders Loop:} Transition from the current offline classical preprocessor to an online iterative loop where the QPU directly executes the Master Problem QUBO solver in each iteration, executing the full Samplomatic error-mitigated pipeline (PEC + SLC + PNA) on IBM Heron.
  \item \textbf{Optimality Cuts:} Implement dual-based optimality cuts to bound and optimize objective function values, moving beyond feasibility-only cuts.
\end{itemize}

% =====================================================
% FIGURES PAGE
% =====================================================
\newpage
\section*{Figures}
\begin{figure}[htbp]
\centering
\includegraphics[width=0.85\textwidth]{hybrid_architecture.png}
\caption{Hybrid Benders-QAMOO architecture: classical Master Problem (MP) and AC power-flow Subproblem (SP) partition, with quantum QAOA sampling the penalized objective graphs.}
\end{figure}

\begin{figure}[htbp]
\centering
\includegraphics[width=0.7\textwidth,height=0.4\textheight,keepaspectratio]{benders_qamoo.png}
\caption{Benders-QAMOO workflow: iterative feasibility-cut generation from AC power-flow violations, compiled into QUBO penalty terms for the quantum variational sampler.}
\end{figure}

\newpage
\bibliographystyle{unsrt}
\bibliography{references}

\end{document}
